\documentclass{article}
\usepackage[T1]{fontenc}

\usepackage[a4paper, margin=2.5cm]{geometry}
\usepackage[polish, bidi=basic, provide=*]{babel}
\babelprovide[import, onchar=ids fonts]{polish}

\usepackage{amsmath}
\usepackage{booktabs}
\usepackage{hyperref}

\title{Raport z badania porównawczego systemów syntezy mowy (TTS)}
\author{TODO: imiona, nazwiska i numery indeksów}
\date{\today}

\begin{document}

\maketitle

\section{Wstęp i cel badania}
Celem niniejszego badania jest ocena przydatności trzech wybranych systemów syntezy mowy
polskiej (TTS) -- Edge-TTS, Piper TTS oraz ElevenLabs -- do obsługi projektowanego interfejsu
konwersacyjnego.

Systemy zostały ocenione pod kątem poprawności wymowy, zrozumiałości, naturalności brzmienia,
intonacji i prozodii, płynności, stabilności syntezy, szybkości generowania oraz przydatności
w dialogu z użytkownikiem.

\section{Wybrane systemy TTS}
\begin{itemize}
    \item \textbf{Edge-TTS} -- darmowy, działa w chmurze (Microsoft), głos: TODO.
    \item \textbf{Piper TTS} -- lokalny, open-source, działa offline, model/głos: TODO.
    \item \textbf{ElevenLabs} -- płatny/limitowany darmowy tier, chmura, głos: TODO.
\end{itemize}
TODO: opisać ustalone porównywalne parametry (szybkość mówienia, sposób generowania).

\section{Metodologia i zbiór testowy}
\subsection{Zbiór danych}
Badanie przeprowadzono na wspólnym zbiorze testowym przygotowanym przez grupę projektową,
składającym się z TODO tekstów od TODO osób (5 tekstów na osobę), podzielonych na 5 kategorii:
\begin{itemize}
    \item Krótkie komunikaty informacyjne,
    \item Pytania i odpowiedzi,
    \item Wypowiedzi zawierające liczby, daty, godziny i skróty,
    \item Zdania złożone i dłuższe wypowiedzi,
    \item Wypowiedzi zawierające nazwy własne i wyrazy obcojęzyczne.
\end{itemize}
Do dalszej, pogłębionej analizy wybrano 15 tekstów własnych grupy oraz 30 tekstów ze
wspólnej bazy (liczby z brief'u $\times 1{,}5$, proporcjonalnie do 3-osobowego składu grupy;
TODO: lista wybranych tekstów).

\subsection{Generowanie i anonimizacja nagrań}
Dla każdego tekstu wygenerowano nagrania wszystkimi trzema systemami
(\texttt{scripts/generate\_speech.py}), a następnie nagrania zanonimizowano względem systemu
źródłowego (\texttt{scripts/anonymize\_audio.py}), aby ograniczyć wpływ oczekiwań oceniających.

\subsection{Kryteria oceny}
Kryteria 2--6 i 8 oceniono w 5-punktowej skali Likerta (1 -- bardzo słaba, 5 -- bardzo dobra).
Kryterium 1 (poprawność wymowy) oceniono poprzez zliczanie błędów w kategoriach: liczby/daty/
godziny, skróty, nazwy własne, wyrazy obce, akcentowanie. Kryterium 7 (szybkość generowania)
zmierzono jako czas od przekazania tekstu do syntezatora do wygenerowania dźwięku.

\section{Wyniki badań}

\subsection{Poprawność wymowy (kryterium 1)}
TODO: tabela najczęstszych błędów per system (\texttt{results/bledy\_wymowy-NNN.csv}).

\subsection{Oceny subiektywne (kryteria 2--6, 8)}
TODO: tabela średnich ocen per system i kryterium, ogółem i per kategoria
(\texttt{results/podsumowanie\_tts-NNN.csv}).

\subsection{Szybkość generowania (kryterium 7)}
TODO: tabela średniego czasu generowania i znaków/sekundę per system
(\texttt{results/latency\_tts-NNN.csv}).

\subsection{Ocena końcowa}
TODO: wynik \texttt{scripts/compute\_final\_score.py} -- UWAGA: wagi domyślne w tym skrypcie
są placeholderem (1/3 na komponent), nie wynikiem ustalonym z wymaganiami ćwiczenia.

\section{Analiza wyników}
TODO:
\begin{enumerate}
    \item Porównanie wyników dla różnych kategorii wypowiedzi.
    \item Najczęściej występujące błędy syntezy.
    \item Porównanie czasu generowania mowy przez poszczególne systemy.
    \item Który system najlepiej sprawdza się w komunikacji z użytkownikiem.
\end{enumerate}

\section{Podsumowanie i wnioski końcowe}
TODO.

\end{document}
